\scanpage{14}
(par exemple courbe elliptique) est tel que ${}_pA_K$ soit un groupe constant, alors est-il vrai que la trace $\phi_E$ soit \deleted{nulle}\inserted{égale à $n=\dim A_K$} sur les \deleted{\unsure{premtrs}} éléments premiers à $p$ de $I$ ? (Seul le cas de réduction non semi-stable de $A$ offre un problème; donc il faudrait ici trouver un exemple où le critère de réduction semi-stable de Raynaud est en défaut pour $\ell=p$ !).

Dans le cas d'une courbe elliptique, et en car. rés. $\ne 2,3$, prenant pour $K\inserted{'}$ la plus petite extension sur laquelle on a réduction semi-stable, on trouve que $G$ est un sous-groupe de $\operatorname{Aut}(A_{s'})$, lequel \inserted{est cyclique} a $2,4$ ou $6$ éléments donc est premier à $p$, et la question revient à celle de savoir si $\phi_E$ est de valeur\deleted{\unsure{s}} constante égale à $2$, \deleted{si}\inserted{lorsque} ${}_pA$ est constant\deleted{\unsure{,}} (et si $K$ contient les racines p.èmes de l'unité\inserted{(?)}; il suffit de regarder le cas où la courbe elliptique réduite, après extension $K\longrightarrow K'$, est d'invariant de Hasse nul\inserted{)}\deleted{\unsure{.}}
\margin{\uline{faux \unsure{par}} \uline{Tate}}
\mspage{13}
